\section{Accuracy-Matched FP8 Control}
\label{app:accuracy-matched-control}

The principal NV/MX result is a speed--accuracy trade-off, so a separate
control asks what happens when the TK route is required to match HAO's
reported NV/FP8 cosine. Table~\ref{tab:accuracy-matched-control} uses the
B1/S32768/H24/D128 B300 record. The HAO rows are reconstructed from published
TFLOP/s and are cross-run context \cite{hao2026fp4flashattention}; HAO does
not publish relative-$L_2$.

\begin{table}[ht]
\centering
\caption{Accuracy-matched B300 control. Superscript p marks HAO-published
cross-run values.}
\label{tab:accuracy-matched-control}
\scriptsize
\begin{adjustbox}{max width=\columnwidth}
\begin{tabular}{lrrrr}
\toprule
Provider & Time (ms) & TFLOP/s & Cosine & Relative-$L_2$ \\
\midrule
TK NV/MX fast & 4.481 & 2945 & 0.9429 & 0.3389 \\
TK NV/FP8 optimized & 7.584 & 1740 & 0.9573 & 0.2913 \\
TK NV/FP8 exact & 8.854 & 1490 & 0.9897 & 0.1433 \\
HAO NV/FP8\textsuperscript{p} & 4.929 & 2677 & 0.9899 & -- \\
HAO BF16\textsuperscript{p} & 8.607 & 1533 & 1.0000 & -- \\
\bottomrule
%
\end{tabular}
\end{adjustbox}
\end{table}

The exact TK NV/FP8 route reaches 0.9897 cosine, effectively matching HAO's
published 0.9899 at the displayed precision. It takes 8.854 ms and reaches
1490 TFLOP/s, compared with HAO's 4.929 ms and 2677 TFLOP/s. It also slightly
trails HAO's published BF16 throughput. The intermediate optimized TK NV/FP8
route improves speed but reaches only 0.9573 cosine. This experiment separates
two conclusions: the HAO-derived pipeline is structurally viable, while the
large speed of our retained NV/MX mode comes from making the exposed P path
cheaper and accepting lower operator fidelity. Matching HAO's accuracy with
the current exact TK FP8 path gives back that advantage.
